% Agent 04: Pages 357--360 (PDF pp.~7--8 + top of p.~9)
% Coverage: Section 2.2 Free-Bound Radiation (end, from Eq.~2.2.6),
%           Section 2.3 Thermal Cyclotron and Synchrotron Radiation (through Eq.~2.3.10)
% Equations (2.2.6)--(2.2.9), (2.3.1)--(2.3.10)

%%% ====================================================================
%%% SECTION 2.2  Free-Bound Radiation (continued from agent\_03)
%%% ====================================================================

The total emissivity, $\varepsilon = \int \varepsilon_\nu \, d\nu$, due to a given free-bound transition is
\begin{align}
\varepsilon &= 8\sqrt{3}\,c^2
\left(\frac{C_A f_i Z^2}{A n_0^2}\right)
\left(\frac{\alpha^2 m_e c^2}{m_p^2}\right)
\left(\frac{2\pi m_e c^2}{3kT}\right)^{1/2}
\rho_0
\left(\frac{C_A f_i Z^2}{A n_0^2}\right)
\left(\frac{\rho_0}{A}\right)
\;\frac{\text{ergs}}{\text{g\,sec}}\;
T_{\rm K}^{-1/2}\;\bar{G}_{bf}
\notag \\
&= \left(4.2 \times 10^{26}
\left|\frac{C_A f_i Z^2}{A n_0^2}\right|
\;\frac{\text{ergs}}{\text{g\,sec}}\right)
\left(\frac{\rho_0}{A}\right)\,
T_{\rm K}^{-1/2}\;\bar{G}_{bf}\,,
\tag{2.2.6}
\end{align}
where $\bar{G}_{bf}$, the mean Gaunt factor, depends only on temperature and on the
bound state, and is approximately equal to one
\begin{equation}
\bar{G}_{bf} = \int_0^\infty G_{bf}(h\nu - E_i = xkT;\; n_i, q_i, \ldots)\; e^{-x}\,dx \simeq 1.
\tag{2.2.7}
\end{equation}

The total emissivity~(2.2.6) due to a given free-bound transition, compared
to the total emissivity~(2.1.29) of free-free radiation from the same type of ion,
is
\begin{equation}
\frac{\varepsilon_{fb}}{\varepsilon_{ff}}
\simeq \frac{Z^2}{n_i^4}\left(\frac{8 \times 10^5}{T}\right).
\tag{2.2.8}
\end{equation}

Ginzburg~\cite{Ginzburg1967} states that for astrophysical plasmas the ratio of total free-bound
radiation to total free-free radiation (summed over all states and all ions)
generally is less than
\begin{equation}
\frac{\varepsilon_{fb\;\text{total}}}{\varepsilon_{ff\;\text{total}}}
< \left(\frac{8 \times 10^5\;\text{K}}{T}\right).
\tag{2.2.9}
\end{equation}

%%% ====================================================================
%%% SECTION 2.3  Thermal Cyclotron and Synchrotron Radiation
%%% ====================================================================

\subsection{Thermal Cyclotron and Synchrotron Radiation}
\label{sec:2.3}

\textit{Basic references}\quad Jackson~\cite{Jackson1962};\
Ginzburg~\cite{Ginzburg1967}.

\bigskip
\noindent\textit{Power radiated by a single electron}

\medskip
Consider an electron of speed~$v$ and total mass-energy~$\gamma m_e c^2$,
\begin{equation}
\gamma \equiv (1 - v^2)^{-1/2},
\tag{2.3.1}
\end{equation}
spiralling in a magnetic field of strength~$B$. The Lorentz 4-acceleration of the
electron is
\begin{equation}
a^0 = 0, \qquad
\mathbf{a} = (e / m_e c)\,\gamma\,\mathbf{v} \times \mathbf{B}.
\tag{2.3.2}
\end{equation}
Consequently, the total power radiated---as obtained from the standard
Lorentz-invariant equation,
\begin{equation}
\frac{dE}{dt} = \frac{2e^2}{3c^3}\,a^2
\tag{2.3.3}
\end{equation}

\noindent---is
\begin{equation}
\frac{dE}{dt} = \frac{2\,r_0^2}{3\,c}\,(\gamma v_\perp)^2\,B^2.
\tag{2.3.4}
\end{equation}
Here $v_\perp$ is the component of the electron velocity perpendicular to the magnetic
field. If the electron is nonrelativistic, $v \ll c$, this radiation is called ``cyclotron
radiation''. If the electron is ultrarelativistic, $\gamma \gg 1$, it is called ``synchrotron
radiation''.

\bigskip
\noindent\textit{Cyclotron radiation from a nonrelativistic plasma}

\medskip
Consider a plasma with temperature
\begin{equation}
kT \ll m_e c^2, \quad \text{i.e.}\ T \lesssim 6 \times 10^9\;\text{K},
\tag{2.3.5}
\end{equation}
and with a magnetic field of strength~$B$. As in the nonrelativistic case, radiation
from protons spiralling
in the magnetic field is smaller by a factor $(m_e/m_p)^3 \simeq 10^{-10}$ than that from
electrons, since
\begin{equation}
\frac{dE}{dt} \propto \frac{v^2}{m^2} \propto \frac{1}{m^3}\,.
\tag{2.3.6}
\end{equation}
(Recall: in thermal equilibrium the mean kinetic energies of protons and electrons
are the same.) Hence, we can ignore protons and other ions when calculating
cyclotron radiation. Let $f_e$ be the number of unbound electrons per baryon in
the plasma. Then the total emissivity (ergs per second per gram of plasma) is
\begin{equation}
\varepsilon = f_e \left\langle \frac{dE}{dt} \right\rangle
\bigg/\!\left(\frac{1}{m_p}\right),
\tag{2.3.7}
\end{equation}
where $\langle\;\rangle$ denotes an average over the Maxwell-Boltzmann velocity distribution
of the electrons. Since $\gamma = 1$, since 2 of the 3 spatial directions are orthogonal
to~$B$, and since the mean kinetic energy per electron is $\tfrac{3}{2}kT$, we have
\begin{equation}
\langle \gamma v_\perp^2 \rangle
= \frac{2}{3}\,\frac{2\,\bar{E}}{m_e}
= \frac{2}{3}\,\frac{3kT}{m_e}
= \frac{2kT}{m_e}\,.
\tag{2.3.8}
\end{equation}
Combining this with equations~(2.3.4) and~(2.3.7), we obtain for the total
emissivity of the plasma
\begin{align}
\varepsilon &= \frac{4}{3}\left(\frac{f_e\,r_0^2}{m_p}\right)
\left(\frac{kT}{m_e c^2}\right) B^2
\tag{2.3.9}\\
&= (0.32\;\text{ergs/g\,sec})\;f_e\;T_{\rm K}\;B_G^2\,,
\notag
\end{align}
where $T_{\rm K}$ is temperature in~${}^\circ$K and $B_G$ is magnetic field in Gauss.

Each electron emits its radiation monochromatically, with the cyclotron
frequency (orbital frequency of electron's spiral motion)
\begin{equation}
\nu_{\text{cyc}} = \frac{eB}{2\pi m_e c}
= (2.79\;\text{MHz})\;B_G\,.
\tag{2.3.10}
\end{equation}

Thus, if the magnetic field is uniform, the radiation will be monochromatic
with this frequency. But in Nature the magnetic field will always be inhomo\-geneous,
and the spectrum will show a spread proportional to the spread in~$B^2$.
(Of course, there are always other sources of spread in the spectrum---e.g., doppler
shifts and relativistic generation of harmonics.)
